\documentclass{standalone}
\usepackage{circuitikz}
\usetikzlibrary{calc}
\definecolor{circuitnoteink}{HTML}{000000}
\definecolor{circuitnotemark}{HTML}{FE00FE}
\begin{document}
\begin{circuitikz}[american, line width=0.8pt]
\ctikzset{bipoles/length=1.2cm}
\ctikzset{grounds/scale=1.36}
\coordinate (b3) at (4,-2);
\coordinate (b6) at (10,-2);
\coordinate (b9) at (16,-2);
\coordinate (b12) at (22,-2);
\draw (b3) node[ocirc]{} node[above left]{IN}; % line 3
\node[ground] at (b6) {}; % line 4
\node[vcc] at (b9) {VCC}; % line 5
\node[vee] at (b12) {VEE}; % line 6
\draw[circuitnoteink] (1,-0.2) -- (25,-0.2); % line 8
\draw[circuitnoteink] (1,-4.2) -- (25,-4.2); % line 9
\draw[circuitnoteink] (1,-0.2) -- (1,-4.2); % line 10
\draw[circuitnoteink] (7,-0.2) -- (7,-4.2); % line 11
\draw[circuitnoteink] (13,-0.2) -- (13,-4.2); % line 12
\draw[circuitnoteink] (19,-0.2) -- (19,-4.2); % line 13
\draw[circuitnoteink] (25,-0.2) -- (25,-4.2); % line 14
\path (3.175,-1.393) rectangle (4.826,-0.207); % line 15
\node[anchor=west, inner sep=0, circuitnotemark, font=\large] at (4,-0.8) {X}; % line 15
\path (2.603,-3.993) rectangle (5.397,-2.807); % line 16
\node[anchor=west, inner sep=0, circuitnotemark, font=\large] at (4,-3.4) {X}; % line 16
\path (8.508,-1.393) rectangle (11.492,-0.207); % line 17
\node[anchor=west, inner sep=0, circuitnotemark, font=\large] at (10,-0.8) {X}; % line 17
\path (8.349,-3.993) rectangle (11.651,-2.807); % line 18
\node[anchor=west, inner sep=0, circuitnotemark, font=\large] at (10,-3.4) {X}; % line 18
\path (14,-1.393) rectangle (18,-0.207); % line 19
\node[anchor=west, inner sep=0, circuitnotemark, font=\large] at (16,-0.8) {X}; % line 19
\path (14.603,-3.993) rectangle (17.397,-2.807); % line 20
\node[anchor=west, inner sep=0, circuitnotemark, font=\large] at (16,-3.4) {X}; % line 20
\path (20,-1.393) rectangle (24,-0.207); % line 21
\node[anchor=west, inner sep=0, circuitnotemark, font=\large] at (22,-0.8) {X}; % line 21
\path (20.476,-3.993) rectangle (23.524,-2.807); % line 22
\node[anchor=west, inner sep=0, circuitnotemark, font=\large] at (22,-3.4) {X}; % line 22
\node[anchor=south west, inner sep=2pt, circuitnotemark, font=\LARGE\bfseries] (circuittitle) at (current bounding box.north west) {X};
\path (circuittitle.south west) rectangle ++(6.736,0.853);
\end{circuitikz}
\end{document}
